\begin{table}[htbp]
\centering
\begin{threeparttable}
\caption{Falsification and Placebo Tests}
\label{tab:placebo_tests}
{\small
\begin{tabular}{lcc}
\toprule
 & Estimate & $p$-value \\
\midrule
\multicolumn{3}{l}{\textit{Panel A: Permutation Placebo IV}} \\
\addlinespace
Mean coefficient & 14.845 & \\
2.5th percentile & $-129.691$ & \\
97.5th percentile & 76.854 & \\
\addlinespace
\multicolumn{3}{l}{\textit{Panel B: Placebo Outcomes and Anticipation}} \\
\addlinespace
Employment expectation index & 986.714 & 0.783 \\
Lead revision (anticipation test) & $-15.333$ & 0.846 \\
\bottomrule
\end{tabular}
}
\begin{tablenotes}[flushleft]
\footnotesize
\item \textit{Notes:} Permutation placebo reassigns the instrument across quarters 300 times. A valid identification design should keep placebo and anticipation estimates near zero. The employment expectation index is a PBoC diffusion index bounded between approximately 44 and 54 (SD $\approx 2.3$ index points, excluding two apparent data-entry outliers in 2014Q3 and 2015Q1). The coefficient of 986.714 is a 2SLS point estimate produced by a weak first stage (Wald $F = 1.16$); weak-instrument bias inflates IV coefficients toward the ratio of reduced-form noise to a near-zero first stage. The associated $p$-value of 0.783 confirms the estimate is statistically indistinguishable from zero, consistent with the placebo null.
\end{tablenotes}
\end{threeparttable}
\end{table}
